\documentclass[11pt]{article}

\usepackage[utf8]{inputenc}
\usepackage[T1]{fontenc}
\usepackage{geometry}
\usepackage{amsmath}
\usepackage{amssymb}
\usepackage{booktabs}
\usepackage{hyperref}
\usepackage{xurl}
\usepackage{xcolor}
\usepackage{listings}
\usepackage{enumitem}
\usepackage{parskip}
\usepackage{graphicx}
\graphicspath{{../images/}}

\geometry{margin=1in}

\lstset{
  language=Java,
  basicstyle=\ttfamily\small,
  keywordstyle=\color{blue},
  commentstyle=\color{gray},
  stringstyle=\color{teal},
  showstringspaces=false,
  columns=fullflexible,
  frame=single,
  framerule=0.2pt,
  breaklines=true
}

\setlist[itemize]{topsep=0.3em,itemsep=0.2em}
\setlist[enumerate]{topsep=0.3em,itemsep=0.2em}

% Simple per-section source line.
\newcommand{\notesources}[1]{\par\noindent{\small\color{gray}\textit{Sources:} \sloppy #1\par}}

% Unit-wise source strings for this subject (used by slides/notes).
% Path is relative to the lecture `latex/` folder (the compilation CWD).
% OOPS with Java (CSEG1044) — unit-wise sources
%
% These are short, student-facing reference strings intended to be used with:
%   - \slidesources{...}  (in beamer slides)
%   - \notesources{...}   (in lecture notes)
%
% Style inspiration: other subjects in My-Subjects/ use semicolon-separated sources
% and include an "accessed" date for web documentation.

\newcommand{\OOPSUnitISources}{Schildt, \textit{Java: A Beginner's Guide} (10e), McGraw-Hill (Java fundamentals + OOP basics).; Oracle Java Tutorials: Getting Started + Language Basics + Classes and Objects (accessed \today).; Java SE API docs (e.g., \texttt{String}, \texttt{Scanner}) (accessed \today).}

\newcommand{\OOPSUnitIISources}{Schildt, \textit{Java: The Complete Reference} (12e), McGraw-Hill (classes, inheritance, packages, interfaces).; Bloch, \textit{Effective Java} (3e), Addison-Wesley, 2018 (design choices; interfaces vs abstract classes).; Oracle Java Tutorials: Classes and Objects + Interfaces + Packages (accessed \today).; Java SE API docs (e.g., \texttt{Comparable}, \texttt{Runnable}) (accessed \today).}

\newcommand{\OOPSUnitIIISources}{Schildt, \textit{Java: The Complete Reference} (12e) (nested/inner classes, exceptions, threads, I/O).; Oracle Java Tutorials: Nested Classes + Exceptions + Concurrency + Basic I/O (accessed \today).; Java SE API docs (e.g., \texttt{Thread}, \texttt{AutoCloseable}) (accessed \today).}

\newcommand{\OOPSUnitIVSources}{Schildt, \textit{Java: The Complete Reference} (12e) (generics, lambdas, Swing, JDBC).; Bloch, \textit{Effective Java} (3e) (generics + lambdas).; Oracle Java Tutorials: Generics + Lambda Expressions + Swing + JDBC Basics (accessed \today).; Java SE API docs (e.g., \texttt{java.util.function}, \texttt{javax.swing}, \texttt{java.sql}) (accessed \today).}

\newcommand{\OOPSUnitVSources}{Schildt, \textit{Java: The Complete Reference} (12e) (Collections Framework + wrapper classes).; Oracle Java docs: Collections Framework overview (accessed \today).; Java SE API docs (\texttt{java.util} collections; wrapper classes in \texttt{java.lang}) (accessed \today).}

\newcommand{\OOPSUnitVISources}{Git SCM: \textit{Pro Git} (2e) (version control workflow), accessed \today.; GitHub Docs (repositories, issues, pull requests), accessed \today.; Oracle Java Tutorials: Swing + JDBC (accessed \today).; JUnit 5 User Guide (accessed \today).; Apache Maven Project: Getting Started (accessed \today).; Gradle User Manual: Getting Started (accessed \today).; UML class diagrams: UML 2.x overview/spec (accessed \today).}

\newcommand{\OOPSMiscWhyInterfacesSources}{Schildt, \textit{Java: The Complete Reference} (12e) (interfaces).; Bloch, \textit{Effective Java} (3e) (prefer interfaces where appropriate).; Oracle Java Tutorials: Interfaces (accessed \today).; Java SE API docs (e.g., \texttt{Comparable}, \texttt{Runnable}) (accessed \today).}



\title{Object Oriented Programming (CSEG1044)\\Unit 05 -- Lecture 02 Notes\\Set/SortedSet + Map/SortedMap}
\author{Tofik Ali}
\date{\today}

\begin{document}
\maketitle
\tableofcontents

\section{Lecture Overview}
This lecture focuses on \textbf{uniqueness} (Set) and \textbf{key-value lookup} (Map).
These structures are essential for problems like:
\begin{itemize}
  \item removing duplicates,
  \item membership checks (``is this item present?''),
  \item counting and grouping (frequency tables),
  \item fast search by key (roll number $\rightarrow$ student details).
\end{itemize}
\notesources{\OOPSUnitVSources}

\section{Syllabus Alignment (Unit 05)}
\textbf{Unit 05:} Set/SortedSet; Map/SortedMap; equality + hashing

\section{Learning Outcomes}
After this lecture, you should be able to:
\begin{itemize}
  \item Choose Set vs List; \texttt{HashSet} vs \texttt{TreeSet}.
  \item Explain \texttt{equals}/\texttt{hashCode} and why they matter in hash-based collections.
  \item Use \texttt{HashMap}/\texttt{TreeMap} for counting/grouping and sorted keys.
\end{itemize}

\section{Key Terms (Quick)}
\begin{itemize}
  \item \textbf{Set}: collection with no duplicates.
  \item \textbf{SortedSet}: set that maintains sorted order (e.g., \texttt{TreeSet}).
  \item \textbf{Hashing}: converting an object to an int hash value (\texttt{hashCode}).
  \item \textbf{Map}: key-value association.
  \item \textbf{SortedMap}: map with sorted keys (e.g., \texttt{TreeMap}).
  \item \textbf{Contract}: if \texttt{a.equals(b)} then \texttt{a.hashCode() == b.hashCode()} must hold.
\end{itemize}

\section{Detailed Notes}
\subsection{Set vs List; HashSet vs TreeSet}
\textbf{List:}
\begin{itemize}
  \item ordered, duplicates allowed, index-based access.
\end{itemize}
\textbf{Set:}
\begin{itemize}
  \item no duplicates, typically no index access.
  \item ideal for membership: ``contains?'' and deduplication.
\end{itemize}

\textbf{HashSet:}
\begin{itemize}
  \item uses hashing; order is not guaranteed.
  \item fast average-case add/contains/remove.
\end{itemize}

\textbf{TreeSet:}
\begin{itemize}
  \item keeps elements sorted (needs \texttt{Comparable} or a \texttt{Comparator}).
  \item slightly slower but provides sorted iteration.
\end{itemize}

\subsection{Equality and hashing basics (must know)}
Hash-based collections (\texttt{HashSet}, \texttt{HashMap}) use:
\begin{itemize}
  \item \texttt{hashCode()} to choose a bucket,
  \item \texttt{equals()} to confirm identity inside that bucket.
\end{itemize}

\textbf{Contract:}
\begin{itemize}
  \item If two objects are equal (\texttt{equals} returns true), their hash codes must be equal.
  \item If you override \texttt{equals}, you should also override \texttt{hashCode}.
\end{itemize}

\textbf{Common mistakes:}
\begin{itemize}
  \item overriding only one of \texttt{equals}/\texttt{hashCode},
  \item using mutable fields in keys (changing key after insertion breaks map/set behavior),
  \item using \texttt{==} instead of \texttt{equals} for strings.
\end{itemize}

\subsection{Map basics; HashMap vs TreeMap}
\textbf{Map} stores pairs: key $\rightarrow$ value.
Common operations:
\begin{itemize}
  \item \texttt{put(k, v)}, \texttt{get(k)}, \texttt{containsKey(k)}, \texttt{remove(k)}.
\end{itemize}

\textbf{HashMap:} no ordering, fast average-case lookup.
\textbf{TreeMap:} keys sorted, good when you need ordered iteration or range queries.
\notesources{\OOPSUnitVSources}

\section{Worked Example: set + map counting}
\begin{lstlisting}
import java.util.*;

public class Demo {
    public static void main(String[] args) {
        // Dedup using Set
        List<String> names = Arrays.asList("Asha", "Ravi", "Asha", "Neha");
        Set<String> unique = new HashSet<>(names);
        System.out.println("Unique: " + unique);

        // Counting with Map (frequency table)
        String[] words = {"a", "b", "a", "c", "b", "a"};
        Map<String, Integer> freq = new HashMap<>();
        for (String w : words) {
            freq.put(w, freq.getOrDefault(w, 0) + 1);
        }
        System.out.println("Freq: " + freq);

        // Sorted keys using TreeMap
        Map<String, Integer> sorted = new TreeMap<>(freq);
        System.out.println("Sorted freq: " + sorted);
    }
}
\end{lstlisting}

\section{In-Class Exercises (with brief solutions)}
\begin{enumerate}
  \item When should you prefer a Set over a List?\par
  \textbf{Answer:} when duplicates are not allowed or you need fast membership checks.

  \item Which collection keeps elements sorted: \texttt{HashSet} or \texttt{TreeSet}?\par
  \textbf{Answer:} \texttt{TreeSet}.

  \item If you override \texttt{equals}, what else should you override?\par
  \textbf{Answer:} \texttt{hashCode}.
\end{enumerate}

\section{Self-Study Practice (no solutions)}
\begin{enumerate}
  \item Count character frequency in a string using \texttt{Map<Character, Integer>}.
  \item Store student details in a map: rollNo $\rightarrow$ Student object.
  \item Create a \texttt{TreeSet} of integers and print them in sorted order.
\end{enumerate}

\section{Checkpoint Questions (with sample answers)}
\textbf{Q1.} Why are \texttt{equals} and \texttt{hashCode} important for \texttt{HashSet}?\par
\textbf{Sample answer:} \texttt{HashSet} uses \texttt{hashCode} to locate a bucket and \texttt{equals} to check if an element already exists. Incorrect implementations can cause duplicates or missing lookups.

\vspace{0.6em}
\textbf{Q2.} When do you choose \texttt{TreeMap} over \texttt{HashMap}?\par
\textbf{Sample answer:} Choose \texttt{TreeMap} when you need keys in sorted order or want ordered iteration/range queries. Use \texttt{HashMap} when ordering is not needed and you want faster average lookup.

\end{document}
